% GENERATED FILE. Edit canonical lessons, then run npm run generate:books.
% canonical-source-hash: fnv1a64:a252093a25578542
% canonical-lessons: FA-C81-shanbe, FA-C81-yekshanbe, FA-C81-doshanbe, FA-C81-chaharshanbe, FA-C81-panjshanbe

\chapter{Saturday, Sunday, Monday, Wednesday, Thursday}
\label{ch:fa-saturday-sunday-monday-wednesday-thursday}
\hlchaptermodality{81}

\begin{chapteropening}
\textbf{By the end of this chapter:} \emph{I can say Saturday, Sunday, Monday, Wednesday and Thursday.}

Complete the last lesson of chapter 81: I can say Saturday, Sunday, Monday, Wednesday and Thursday.
\end{chapteropening}

\section[\texorpdfstring{\fa{شنبه} (shanbe)}{شنبه (shanbe)}]{\fa{شنبه} (shanbe) — Saturday}

\label{lesson:FA-C81-shanbe}



\begin{quote}
Before the new one: say the Persian for winter, then the Persian for a season.
\end{quote}

\subsection*{You'll want to know: \fa{شنبه}}
\textbf{\fa{شنبه}} — \emph{shanbe} — \textquotedblleft{}Saturday\textquotedblright{}.

\begin{cousinweb}
\emph{shanbe} is the Sabbath word: it is usually traced through Middle Persian \emph{šambad} to the Aramaic and Hebrew word for the Sabbath. Saturday was the Sabbath, and the next weekdays are counted from it.
\end{cousinweb}

\begin{grammarlens}[title={What you've built}]
One more word for this chapter.
\end{grammarlens}

\subsection*{Your turn}
\begin{itemize}
\raggedright
  \item \emph{Say it:} \emph{shanbe}
  \item \emph{Say it:} \emph{shanbe}, once more
  \item \emph{Say it:} say \emph{shanbe}
  \item \emph{Recall:} say the Persian for winter, then the Persian for a season, then say \emph{shanbe} again
\end{itemize}

\begin{tcolorbox}[breakable,colback=teal!4,colframe=teal!35!black,title={Before you move on}]
What does \fa{شنبه} mean? (\textquotedblleft{}Saturday\textquotedblright{}.) Say it once more.
\end{tcolorbox}
\section[\texorpdfstring{\fa{یکشنبه} (yekshanbe)}{یکشنبه (yekshanbe)}]{\fa{یکشنبه} (yekshanbe) — Sunday}

\label{lesson:FA-C81-yekshanbe}



\begin{quote}
Before the new one: say the Persian for a season, then the Persian for Saturday.
\end{quote}

\subsection*{You'll want to know: \fa{یکشنبه}}
\textbf{\fa{یکشنبه}} — \emph{yekshanbe} — \textquotedblleft{}Sunday\textquotedblright{}.

\begin{cousinweb}
\emph{yekshanbe} is \emph{yek}, \textquotedblleft{}one\textquotedblright{}, with \emph{shanbe}, \textquotedblleft{}Saturday\textquotedblright{} (EtymDB): the first day after the Sabbath. \emph{yek} is inherited; \emph{shanbe} is the Sabbath word, usually traced through Middle Persian to Aramaic.
\end{cousinweb}

\begin{grammarlens}[title={What you've built}]
One more word for this chapter.
\end{grammarlens}

\subsection*{Your turn}
\begin{itemize}
\raggedright
  \item \emph{Say it:} \emph{yekshanbe}
  \item \emph{Say it:} \emph{yekshanbe}, once more
  \item \emph{Say it:} say \emph{yekshanbe}
  \item \emph{Recall:} say the Persian for a season, then the Persian for Saturday, then say \emph{yekshanbe} again
\end{itemize}

\begin{tcolorbox}[breakable,colback=teal!4,colframe=teal!35!black,title={Before you move on}]
What does \fa{یکشنبه} mean? (\textquotedblleft{}Sunday\textquotedblright{}.) Say it once more.
\end{tcolorbox}
\section[\texorpdfstring{\fa{دوشنبه} (doshanbe)}{دوشنبه (doshanbe)}]{\fa{دوشنبه} (doshanbe) — Monday}

\label{lesson:FA-C81-doshanbe}



\begin{quote}
Before the new one: say the Persian for Saturday, then the Persian for Sunday.
\end{quote}

\subsection*{You'll want to know: \fa{دوشنبه}}
\textbf{\fa{دوشنبه}} — \emph{doshanbe} — \textquotedblleft{}Monday\textquotedblright{}.

\begin{cousinweb}
\emph{doshanbe} is \emph{do}, \textquotedblleft{}two\textquotedblright{}, with \emph{shanbe}, \textquotedblleft{}Saturday\textquotedblright{} (EtymDB): the second day after the Sabbath. \emph{do} is inherited, from Middle Persian \emph{dō} (EtymDB).
\end{cousinweb}

\begin{grammarlens}[title={What you've built}]
One more word for this chapter.
\end{grammarlens}

\subsection*{Your turn}
\begin{itemize}
\raggedright
  \item \emph{Say it:} \emph{doshanbe}
  \item \emph{Say it:} \emph{doshanbe}, once more
  \item \emph{Say it:} say \emph{doshanbe}
  \item \emph{Recall:} say the Persian for Saturday, then the Persian for Sunday, then say \emph{doshanbe} again
\end{itemize}

\begin{tcolorbox}[breakable,colback=teal!4,colframe=teal!35!black,title={Before you move on}]
What does \fa{دوشنبه} mean? (\textquotedblleft{}Monday\textquotedblright{}.) Say it once more.
\end{tcolorbox}
\section[\texorpdfstring{\fa{چهارشنبه} (chahârshanbe)}{چهارشنبه (chahârshanbe)}]{\fa{چهارشنبه} (chahârshanbe) — Wednesday}

\label{lesson:FA-C81-chaharshanbe}



\begin{quote}
Before the new one: say the Persian for Sunday, then the Persian for Monday.
\end{quote}

\subsection*{You'll want to know: \fa{چهارشنبه}}
\textbf{\fa{چهارشنبه}} — \emph{chahârshanbe} — \textquotedblleft{}Wednesday\textquotedblright{}.

\begin{cousinweb}
\emph{chahârshanbe} is \emph{chahâr}, \textquotedblleft{}four\textquotedblright{}, with \emph{shanbe}, \textquotedblleft{}Saturday\textquotedblright{} (EtymDB): the fourth day after the Sabbath. \emph{chahâr} is inherited from Middle Persian (EtymDB).
\end{cousinweb}

\begin{grammarlens}[title={What you've built}]
One more word for this chapter.
\end{grammarlens}

\subsection*{Your turn}
\begin{itemize}
\raggedright
  \item \emph{Say it:} \emph{chahârshanbe}
  \item \emph{Say it:} \emph{chahârshanbe}, once more
  \item \emph{Say it:} say \emph{chahârshanbe}
  \item \emph{Recall:} say the Persian for Sunday, then the Persian for Monday, then say \emph{chahârshanbe} again
\end{itemize}

\begin{tcolorbox}[breakable,colback=teal!4,colframe=teal!35!black,title={Before you move on}]
What does \fa{چهارشنبه} mean? (\textquotedblleft{}Wednesday\textquotedblright{}.) Say it once more.
\end{tcolorbox}
\section[\texorpdfstring{\fa{پنجشنبه} (panjshanbe)}{پنجشنبه (panjshanbe)}]{\fa{پنجشنبه} (panjshanbe) — Thursday}

\label{lesson:FA-C81-panjshanbe}



\begin{quote}
Before the new one: say the Persian for Monday, then the Persian for Wednesday.
\end{quote}

\subsection*{You'll want to know: \fa{پنجشنبه}}
\textbf{\fa{پنجشنبه}} — \emph{panjshanbe} — \textquotedblleft{}Thursday\textquotedblright{}.

\begin{cousinweb}
\emph{panjshanbe} is \emph{panj}, \textquotedblleft{}five\textquotedblright{}, with \emph{shanbe}, \textquotedblleft{}Saturday\textquotedblright{} (EtymDB): the fifth day after the Sabbath. \emph{panj} is an inherited Persian number.
\end{cousinweb}

\begin{grammarlens}[title={What you've built}]
One more word for this chapter.
\end{grammarlens}

\subsection*{Your turn}
\begin{itemize}
\raggedright
  \item \emph{Say it:} \emph{panjshanbe}
  \item \emph{Say it:} \emph{panjshanbe}, once more
  \item \emph{Say it:} say \emph{panjshanbe}
  \item \emph{Recall:} say the Persian for Monday, then the Persian for Wednesday, then say \emph{panjshanbe} again
\end{itemize}

\begin{tcolorbox}[breakable,colback=teal!4,colframe=teal!35!black,title={Before you move on}]
What does \fa{پنجشنبه} mean? (\textquotedblleft{}Thursday\textquotedblright{}.) Say it once more.
\end{tcolorbox}
